\documentclass[12pt]{article}
\usepackage[utf8]{inputenc}
\usepackage[T1]{fontenc}
\usepackage{hyperref}
\title{Twi Language Resources for K-12 Curriculum Development}
\author{Compiled by Claude Sonnet 5 (claude-sonnet-5) for Andr\'as Kornai}
\date{}
\begin{document}
\maketitle

\section{Introduction}
Twi (ISO 639-3 twi) is a variety of Akan, spoken mainly in Ghana. It has two literary standards, Asante and Akuapem, and the related Fante has its own orthography. The prompt gave no identification beyond the code, so this report assumes the Akan language of Ghana. Most published courses teach Asante Twi. Only a few resources turned up in the searches.

\section{Primary Educational Resources}
The Cambridge Language Centre list names a comprehensive Twi (Asante) course for non-Twi learners, a basic course, a Pimsleur audio programme and a book of 102 Akan verbs: \url{https://www.langcen.cam.ac.uk/resources/langt/twi.html}. A compiled list of Akan resources mentions the free FSI Akan course (textbook PDF and audio): \url{https://universeofmemory.com/akan-language-resources/}. These are aimed at adult learners, not at school pupils.

\section{Linguistic and Reference Materials}
\begin{itemize}
\item Christaller's 1881 dictionary of the Asante and Fante language, digitised in the University of Ghana repository: \url{https://ugspace.ug.edu.gh/items/4fbc144d-4bbf-43be-9097-308cb7e7d3d6/full}
\item Hippocrene Asante Twi--English dictionary and phrasebook (commercial, about 4,000 entries): \url{https://www.hippocrenebooks.com/store/p44/AsanteTwiDictionaryandPhrasebook.html}
\item The Cambridge list also names a Twi (Asante) learners' reference grammar and a Modern Akan Dictionary: \url{https://www.langcen.cam.ac.uk/resources/langt/twi.html}
\end{itemize}

\section{Religious and Community Resources}
None identified in these searches.

\section{Digital and Audio Resources}
The FSI Akan course audio recordings and the Pimsleur programme are mentioned above. A trilingual Twi--French--English app was mentioned in a result, with few details: \url{https://resulam.com/?p=1282}

\section{Curriculum Development Context}
Twi is a major language with a long written tradition, but these searches found only adult-learner and reference materials. Ghanaian school textbooks, curriculum documents and policy would need a separate search. Developers should choose the dialect standard (Asante, Akuapem or Fante) first.

\end{document}